\chapter*{Preface}
\addcontentsline{toc}{chapter}{Preface}
\markboth{Preface}{Preface}

\section*{Why this book exists}

The Islamic economic tradition contains a class of claims that the
modern academic literature has not been able to engage with. The
believer hears in every Friday sermon that provision is decreed before
birth, that gratitude multiplies wealth, that interest destroys what it
appears to build, that the wealth in some hands carries blessing and in
others does not. The believer is invited, by the homiletic literature,
to trust these claims as elements of the believer's interior life. The
academic Islamic economist, trained in neoclassical methods, treats the
same claims as outside the field: as theology, perhaps, or as folk
psychology, but not as propositions about the economic order that admit
of theoretical treatment.

This silence is not a small matter. The claims in question are
descriptive: they tell us that wealth has properties the standard
formalism does not capture. If the claims are true, the entire edifice
of conventional economic theory --- and the Islamic economics that has
been built on its foundation --- is missing something structural. If
they are false, the believer should be told. What cannot be defended is
the present situation, in which the claims are at once central to the
religious tradition and absent from the theoretical literature that
purports to articulate the tradition's economic content.

This book exists to attempt the third option. It treats the
metaphysical claims as theoretical propositions and asks what
mathematical and philosophical apparatus would be required to test
them, formalize them, and connect them to one another. The attempt is
not new in spirit --- the medieval and early modern Islamic tradition
took these claims with theoretical seriousness, in the work of
al-Ghazali, Ibn al-Arabi, Mulla Sadra, Shah Wali Allah Dehlawi, and
many others --- but it has not been undertaken in the modern academic
register, with the conceptual apparatus that twentieth-century physics,
mathematics, philosophy, and economics now make available. That
apparatus did not exist in the seventeenth century. It does now. The
question is whether it can carry the weight of the metaphysical claims
of a fourteen-century tradition.

The answer this book proposes is: in many cases, yes; in others, with
careful qualification; in still others, not yet, and perhaps not ever.
The work of the book is to make these distinctions precise.

\section*{The position of this volume in a larger project}

This is Volume 0 of a planned multi-volume series. The position is
deliberate. A textbook treatment of Islamic economic theory in the
tradition of Mas-Colell, Whinston, and Green would presuppose
foundational choices that do not currently exist as settled. What are
the right preference axioms when the agent recognizes both worldly and
otherworldly returns? What is the right probability framework when
provision is held to be decreed but striving is held to matter? What is
the right ontology of wealth when the tradition holds that material
goods are configurations of divine self-disclosure rather than
self-subsisting substances? The textbook must answer these questions
before it can begin its formal work. To answer them is the task of
Volume 0.

When this volume is complete, a second volume can take up the formal
microeconomic theory --- preferences, choice, demand, equilibrium ---
on the foundations established here. That work is already partially
drafted. A third volume will treat macroeconomic theory and public
finance with the same rigor; a fourth, Islamic financial theory; a
fifth, comparative economic systems with rigorous case studies of
Pakistan, Malaysia, Iran, and others. Each will build on the
preceding. The series, taken as a whole, is intended to constitute the
first technically defensible articulation of an Islamic economic
science.

This first volume is therefore, in a sense, the volume that justifies
the others. It does not contain the formal economic theory the later
volumes will develop. It contains the philosophical and mathematical
work that the later volumes are entitled to assume.

\section*{Who this book is for}

The book is written for several audiences and the reader should know in
advance which they are.

The first audience is the academic economist who has read mainstream
Islamic economics, found it derivative, and wants to understand whether
there is something distinctively Islamic in the field that is worth
their attention. To this reader the book offers eight worked
treatments of metaphysical claims, each derived with the rigor of
research economics, each yielding empirical predictions or theoretical
implications that do not appear in the standard literature. The
economist is not asked to accept Islamic theology. The economist is
asked to evaluate whether, granted the philosophical foundations the
book defends, the formal results follow. They do.

The second audience is the philosopher --- of physics, of mind, of
religion --- who has noticed that Islamic metaphysics has gone largely
untreated in twentieth-century philosophy and wonders whether the
tradition has anything to contribute to current debates. To this
reader the book offers a sustained argument that the Akbarian
metaphysical tradition (Ibn Arabi and his successors) anticipates,
in its own vocabulary, several of the positions that contemporary
philosophy has been led to by independent argument: analytical
idealism, perspectival probability, the ontological reality of
imagination, the unity of being. The convergence is not coincidence
and it is not concordism. The philosopher is invited to recognize a
sophisticated tradition that twentieth-century philosophy has had no
reason to know about.

The third audience is the Islamic theologian or scholar of Islamic
intellectual history. To this reader the book offers a translation of
the classical claims into the formal vocabulary of contemporary
science --- not a reduction of those claims to formalism, but a
demonstration that the claims have a precise theoretical content that
can be made visible in modern terms. The book is in this sense an
exercise in the genre that Mulla Sadra called \arb{tatbiq} --- the
matching, or putting-into-correspondence, of inherited tradition with
the conceptual instruments of one's own time. It belongs to a
genre the tradition recognizes.

The fourth audience is the believing Muslim, perhaps with technical
training in some field, who has long felt that the metaphysical
claims of the tradition were doing more work than the homiletic
register acknowledged. To this reader the book offers a vindication
not of belief --- belief is not the kind of thing that admits of
mathematical proof --- but of the structural claim that the
tradition has been making true descriptive statements about the world
all along, statements that the modern analytical apparatus is now
capable of recognizing.

\section*{What this book is not}

A few negations. The book is not concordism. Concordism is the
intellectual genre in which a religious tradition's text is matched
against the latest scientific findings to demonstrate that the text
``predicted'' the findings. Concordism is bad science (it cherry-picks)
and bad theology (it makes the religious tradition hostage to the
scientific consensus of the moment). The convergences this book traces
are not predictions. They are structural recognitions: features of the
world that classical materialism failed to see, that modern science has
been forced to acknowledge, and that the Islamic tradition was already
articulating in its own vocabulary because it was working from
different starting premises that turned out to be closer to the truth.

The book is not apologetics. Apologetics defends a religious
tradition against external critique. The work here is descriptive: the
tradition makes certain claims about the world; either those claims are
true or they are not; the question is what they would have to mean and
what they would have to look like in formal terms for them to be true.
Where the answer is ``they are true and here is the demonstration,''
the book says so. Where the answer is ``the formal treatment is not
yet possible,'' the book says that. Where the answer is ``the claim is
real but its operational content remains to be defined by future
research,'' the book says that. The reader gets honest reporting, not
defense.

The book is not a survey of Islamic economics. There are excellent
surveys, including Khan and Mirakhor and others, and the reader who
wants the survey literature will find it elsewhere. The book takes a
specific position, articulates it at length, and defends it against
specific objections. It is a treatise, not a textbook in the survey
sense.

The book is not a finished work. Several of the treatments in
Part III (in subsequent installments) require empirical research that
has not yet been undertaken; the book argues for the design of that
research and reports preliminary evidence where it exists. Several of
the treatments are operational definitions awaiting field test. The
book is honest about this. It is the foundation of a research program,
not the conclusion of one.

\section*{Pedagogical commitments}

A few words on style. The book assumes some familiarity with
intermediate economics --- supply and demand, basic equilibrium
analysis, expected utility theory at the textbook level --- but
assumes no prior exposure to physics, philosophy of mind, or classical
Islamic philosophy. Concepts from these fields are introduced where
needed, defined precisely, and worked through with examples. The
mathematical level is set as high as the argument requires; the
mathematical exposition is set as low as the argument permits. The
reader who has worked through a graduate course in microeconomic
theory will find the treatment elementary in places; the reader who
has not will find it accessible.

Three pedagogical commitments deserve advance notice.

\begin{enumerate}
    \item \textbf{Connectivity.} The chapters of this book are not
    self-contained. Each chapter is connected, by explicit
    cross-reference, to chapters that precede and follow. The reader
    is expected to come to understand the work as a unified whole, not
    as a collection of independent treatments. Forward and backward
    references appear throughout. Where a result of Chapter 2 is used
    in Chapter 11, the use is flagged in Chapter 2 and referenced
    again in Chapter 11. The connectivity is the point.

    \item \textbf{Definitions before deployment.} Every technical term
    is defined before it is used. The first occurrence of a term is
    set in italic and accompanied by either a definition box or a
    footnote pointing to the definition. The reader who picks up the
    book in the middle should be able to find the definition of any
    unfamiliar term within two pages or by consulting the glossary in
    the back matter.

    \item \textbf{Honest signposting.} When the argument is on solid
    ground, the book says so. When it is on contested ground, the
    book says so. When the argument is the author's own and not
    established in the literature, the book says so. The reader is
    never given the false impression that a contested claim is
    settled.
\end{enumerate}

\section*{What the reader will find in this volume}

The present volume contains the front matter and Part I of the larger
work. Part I, comprising four chapters, establishes the methodological
foundations.

\textbf{Chapter 1} states the two failures of contemporary Islamic
economics --- the failure of mechanism and the failure of rigor ---
and motivates the third path that this work pursues.

\textbf{Chapter 2} develops the three-level taxonomy of rigor (metaphor,
formal analogy, causal mechanism) that organizes every treatment in
the book and that, the author argues, should organize every honest
engagement with metaphysical claims.

\textbf{Chapter 3} surveys the textual and traditional foundations:
the Quranic verses, hadith, classical \arb{kalam} debates, the
Akbarian tradition, and the modern scholarship on which the rest of
the work draws.

\textbf{Chapter 4} explains why modern conceptual tools --- from
quantum interpretations to ergodicity economics to analytical
idealism to the Penrose argument from G\"odel --- are needed to make
sense of the Islamic claims, and previews how Part II will develop
each of them.

Subsequent installments will contain Part II (the modern conceptual
apparatus, in five chapters), Part III (the eight worked treatments,
in eight chapters), and Part IV (synthesis and the path forward, in
four chapters).

\section*{Acknowledgements}

The intellectual debts of this work are enormous and the author
acknowledges them at length in the chapters where they bear. A few
require mention in the preface. To the late Toshihiko Izutsu, whose
\textit{Sufism and Taoism} taught a generation of comparative
philosophers what serious comparative metaphysics looks like. To
William Chittick, whose decades of patient translation and exposition
of Ibn Arabi have made the Akbarian tradition newly available. To
Ole Peters, whose work on ergodicity economics has, almost
single-handedly, made the formal case for the Islamic financial
prohibitions intelligible to mainstream economists who could not
hear the religious case for them. To Roger Penrose, whose argument
from G\"odel deserves more attention from non-physicists than it has
received. To Bernardo Kastrup, whose analytical idealism is the
nearest contemporary Western articulation of the Islamic ontology of
\arb{wujud}. To countless teachers of the tradition, in person and
in books, whose patience with the author's slow learning is the only
reason any of this exists at all.

The errors that remain are entirely the author's. The reader who
finds them is asked to write.

\vspace{1em}
\hfill\textit{The author}\\
\hfill April 2026

\cleardoublepage
